\section{Introduction}
\label{sec:introduction}

\subsection{Research Background and Motivation}
\label{sec:background}
Post-harvest agricultural loss represents a critical challenge for global food security, economic stability, and supply chain sustainability \cite{fao2019state}. According to the Food and Agriculture Organization (FAO), approximately 14\% of global food production is lost between harvest and retail distribution \cite{fao2019state}. In regional agrarian economies, post-harvest deterioration across domestic transit corridors is particularly acute. For example, nationwide surveys in India conducted by the Indian Council of Agricultural Research (ICAR-CIPHET) indicate that post-harvest losses of commercial fruits consistently range between 6.70\% and 15.88\% (e.g., 10.39\% in apple, 7.76\% in banana, and 15.88\% in citrus) \cite{jha2015assessment}. These losses disproportionately penalize smallholder farmers and aggregation packhouses, where produce sorting and quality grading have historically relied upon manual, visual inspection. Manual grading is inherently subjective, labour-intensive, prone to operator fatigue, and incapable of scaling to high-throughput sorting facilities \cite{rizzo2023fruit,zhang2024fruit}.

In response to these operational bottlenecks, computer vision and deep learning have emerged as promising non-destructive methodologies for automated fruit freshness assessment, ripeness grading, and surface defect classification \cite{rizzo2023fruit,zhang2024fruit,chen2024classification,gulzar2023fruit}. Convolutional neural networks (CNNs), particularly lightweight architectures such as MobileNetV2 \cite{sandler2018mobilenetv2}, have become central to practical agricultural deployments. Depthwise separable convolutions and inverted residual bottlenecks allow MobileNetV2 to execute on resource-constrained mobile phones and edge sorting devices with parameter budgets under 4 million weights \cite{chen2024classification,sandler2018mobilenetv2}. In controlled academic benchmarks, modern visual classifiers routinely report top-1 accuracy exceeding 95\% across multi-class produce datasets \cite{zhang2024fruit,chen2024classification,gulzar2023fruit}.

\subsection{The Closed-Set Assumption and Its Vulnerabilities in Agricultural Vision}
\label{sec:closed_set_vulnerabilities}
Despite high benchmark accuracy, practical deployment of deep learning models in real-world agricultural environments encounters fundamental reliability barriers. Standard supervised image classifiers operate under a strict \textit{closed-set assumption}: the training, validation, and operational deployment data are presumed to be independent and identically distributed (i.i.d.) over a fixed, finite set of semantic classes, $\mathcal{Y} = \{1, 2, \dots, K\}$ \cite{guo2017calibration,hendrycks2017baseline}. In operational agricultural workflows, however, this assumption is persistently violated. Vision models deployed at packhouses, retail terminals, or farm collection centres confront three distinct, compounding categories of distribution shift:

\begin{enumerate}
\item \textbf{Covariate and Domain Shift:} Images collected across external facilities exhibit pronounced variations in background illumination, sensor optics, specular glare, ambient shadows, and post-harvest handling marks that differ markedly from laboratory training sets \cite{rizzo2023fruit,chen2024classification,abayomi2024fruitq,sultana2022extensive}. Furthermore, biological cultivar diversity---such as regional Asian apple varieties differing in pigmentation and surface russeting from Western commercial cultivars---induces domain shifts that degrade feature representations \cite{sultana2022extensive,bijoy2025fruitvision}.
\item \textbf{Semantic Shift (Unseen Species):} In automated sorting lines, bins, or field crates, models trained exclusively on a small palette of target crops (e.g., apple, banana, and orange) frequently encounter non-target produce species (e.g., mango, guava, jujube, pomegranate, or lime). Under standard closed-set inference, deep networks lack an explicit rejection mechanism; they are forced to assign an unseen species to one of the $K$ known classes, often with uncalibrated, erroneously high confidence \cite{guo2017calibration,hendrycks2017baseline}.
\item \textbf{Treatment-Labelled Condition Shift:} Horticultural commodities in post-harvest trade may undergo preservative treatments intended to retard visual senescence during extended storage or transit. In curated benchmarks such as FruitVision \cite{bijoy2025fruitvision}, images are documented under a specific post-harvest treatment protocol involving immersion in a 10\% formalin and 90\% water solution for 30 minutes, followed by drying and photography 48 hours later. When produce has undergone chemical treatment, superficial visual decay symptoms may be altered, altering the standard visual manifestations of biological freshness.
\end{enumerate}

Under standard softmax classification, neural networks are notoriously overconfident \cite{guo2017calibration}. Logit outputs normalized via the softmax function produce probability distributions that often peak sharply near 1.0 even when the input image is heavily corrupted, semantically anomalous, or completely out-of-distribution \cite{guo2017calibration,hendrycks2017baseline}. In agricultural supply chains, an overconfident false-positive classification---such as misclassifying an unseen non-target species or chemically treated decaying lot as ``Fresh Produce''---presents severe commercial, regulatory, and quality-assurance liabilities.

\subsection{Open-Set Recognition and Out-of-Distribution Detection Foundations}
\label{sec:osr_ood_foundations}
To address closed-set failure modes, computer vision research has developed two closely related paradigms: Open Set Recognition (OSR) \cite{scheirer2013toward,bendale2016towards} and Out-of-Distribution (OOD) Detection \cite{hendrycks2017baseline,yang2024generalized}. 

The concept of open-set recognition was formalized by Scheirer et al. \cite{scheirer2013toward}, who defined the objective as simultaneously classifying known training classes while bounding the open space risk---the risk of incorrectly labeling an unknown, novel class sample as any known class. Bendale and Boult \cite{bendale2016towards} extended this concept to deep neural networks through OpenMax, which estimates the probability of an input belonging to an unknown class by fitting Weibull distributions to activation vectors in the penultimate layer.

Simultaneously, the broader machine learning community formalized generalized out-of-distribution detection \cite{hendrycks2017baseline,yang2024generalized}. Hendrycks and Gimpel \cite{hendrycks2017baseline} established the Maximum Softmax Probability (MSP) baseline, showing that correctly classified in-distribution samples typically exhibit higher maximum posterior probabilities than anomalous inputs. Liang et al. \cite{liang2018enhancing} developed ODIN, which amplifies the separation between in-distribution and out-of-distribution samples using temperature-scaled softmax outputs and input-space gradient perturbations. Lee et al. \cite{lee2018simple} introduced a feature-space Mahalanobis distance metric that models class-conditional Gaussian distributions across intermediate hidden representations. More recently, Liu et al. \cite{liu2020energy} proposed Energy-based OOD detection, utilizing the Helmholtz free energy function computed directly over unnormalized logits. Energy scores possess the theoretical advantage of being linearly aligned with the true input marginal probability density, avoiding the artificial probability saturation inherent in softmax normalization \cite{yang2024generalized,liu2020energy}. Additionally, Vaze et al. \cite{vaze2022open} demonstrated that closed-set classifier capacity and representation quality strongly dictate open-set discrimination performance.

In operational deployment, OOD detectors are typically implemented either as continuous ranking functions or as discrete rejection gates governed by a pre-calibrated threshold $\tau$. When an input sample produces an anomaly score that violates the threshold, the system abstains from classification, triggering manual inspection or secondary non-optical assay.

\subsection{Prior Agricultural OSR/OOD Studies and Specific Research Gaps}
\label{sec:prior_work_gaps}
Open-set recognition and anomaly detection have recently attracted growing attention in agricultural vision. Several recent investigations have examined related challenges:
\begin{itemize}
\item Sattar et al. \cite{sattar2024sensornet} developed SensorNet, combining computer vision, deep learning, and chemical sensors to identify chemical adulteration in fruit specimens.
\item Bijoy et al. \cite{bijoy2025fruitvision} introduced FruitVision, an image dataset containing fresh, rotten, and formalin-mixed fruit images.
\item Wang et al. \cite{wang2023crossdomain} proposed lightweight attention networks with unsupervised domain adaptation for cross-domain fruit classification.
\item Pehlivan and Kara \cite{pehlivan2026lightweight} introduced a lightweight one-shot open-set metric learning framework for food recognition in smart appliances.
\item Jansi and Kavitha \cite{jansi2026peos} developed PEOS-Net, an evidential open-set framework for plant leaf disease recognition.
\item Yan et al. \cite{yan2026pagmsbnet} formulated PAG-MSBNet, using prototype-augmented graphs for open-set crop disease recognition.
\end{itemize}

While these studies highlight the importance of non-standard conditions in agricultural classification, existing research has not systematically evaluated how established post-hoc OOD scoring functions behave when a frozen classifier faces multi-source external distributions. Specifically, four key empirical gaps remain:

\begin{enumerate}
\item \textbf{Gap 1: Over-Optimistic Internal-Split Benchmarks:} The majority of existing fruit-freshness studies report classification and anomaly detection metrics computed exclusively on random train/test splits of a single, curated dataset \cite{zhang2024fruit,chen2024classification,gulzar2023fruit}. Audits reveal that random image partitioning in agricultural benchmarks frequently suffers from near-duplicate and affine leakage between training and testing sets, producing inflated accuracies that fail to generalize to external repositories \cite{abayomi2024fruitq,sultana2022extensive}.
\item \textbf{Gap 2: Cross-Repository Confounding in OOD Benchmarks:} When external evaluations are conducted, researchers commonly evaluate unseen species drawn from a completely different repository than the known-class training cohort. This experimental design introduces severe \textit{cross-repository confounding}: the detector's ability to reject the unseen species is driven largely by differences in camera hardware, studio lighting, and background compression rather than true biological or semantic divergence \cite{yang2024generalized}. Without \textit{source-stratified} benchmarks that evaluate known and unknown produce collected under identical repository conditions, the true semantic discrimination capacity of vision models remains unmeasured.
\item \textbf{Gap 3: Empirical Behavior Under Treatment-Labelled Shifts:} While multi-modal sensing has explored chemical adulteration \cite{sattar2024sensornet}, no prior study has systematically evaluated whether post-hoc OOD detection methods applied to standard RGB classifiers can discriminate between naturally fresh fruit and produce documented under chemical immersion protocols \cite{bijoy2025fruitvision}.
\item \textbf{Gap 4: Retrospective Metric Optimism vs. Operational Threshold Reliability:} Nearly all published OOD literature evaluates algorithms using threshold-free retrospective metrics, specifically the Area Under the Receiver Operating Characteristic curve (AUROC) and the False Positive Rate at 95\% True Positive Rate (FPR@95) \cite{hendrycks2017baseline,yang2024generalized,liu2020energy}. While informative for theoretical comparison, retrospective metrics assume oracle access to the entire test distribution. In real-world sorting, systems must operate prospectively with \textit{fixed thresholds} pre-calibrated on internal validation sets. The degree to which fixed validation thresholds maintain acceptable Unknown False Acceptance Rates (UFAR) without causing severe Known-Class Rejection Rates (KRR) on external produce has not been rigorously quantified.
\end{enumerate}

\subsection{Research Objectives and Formal Research Questions}
\label{sec:research_questions}
This investigation addresses these gaps by conducting an extensive, source-stratified empirical evaluation of out-of-distribution detection in fruit image classification. The study evaluates a frozen MobileNetV2 six-class visual classifier across $N = 13,465$ image records curated from three distinct repositories under the Stage 4.0.3 external evaluation protocol. 

The paper specifically investigates three core research questions:
\begin{itemize}
\item \textbf{Research Question 1 (RQ1 --- External Generalization):} How does the closed-set classification performance (top-1 accuracy, macro-F1, and balanced accuracy) of a lightweight MobileNetV2 model change when evaluated on external repositories containing the same known fruit classes?
\item \textbf{Research Question 2 (RQ2 --- OOD Discrimination):} How effectively do established OOD scoring mechanisms (MSP, Energy, ODIN, Mahalanobis distance, OpenMax) discriminate between in-distribution produce and distribution shifts when evaluated in source-matched cohorts isolating (a) unseen-species semantic shift, (b) treatment-labelled condition shift, and (c) combined species-and-treatment shifts?
\item \textbf{Research Question 3 (RQ3 --- Operational Reliability):} How do fixed validation-calibrated rejection thresholds ($\tau_{95}$) affect out-of-distribution false acceptance (UFAR) and known-class retention (KRR) under external domain and condition shifts?
\end{itemize}

\subsection{Evidence-Supported Research Contributions}
\label{sec:contributions}
The primary contribution of this work is an empirical source-stratified evaluation, rigorous uncertainty audit, and diagnostic failure analysis. The study does not propose a novel standalone OOD algorithm, nor does it claim real-world agricultural supply chain deployment. The four specific, evidence-backed contributions of this study are:

\begin{enumerate}
\item \textbf{Audited Source-Stratified Benchmark Census ($N = 13,465$):} We establish an audited image-level evaluation benchmark comprising 13,465 image records structured across five distinct partitions: Partition A (Primary ID Test, $N=227$), Partition B (External Known-Class Shift, $N=5,350$), Partition C (Treatment-Labelled Condition Shift, $N=1,917$), Partition D1 (Unseen-Species Shift, $N=4,776$), and Partition D2 (Combined Shift, $N=1,195$). We incorporate exact-duplicate handling and perceptual-hash (pHash) sensitivity analysis across Hamming distance radii $r \in \{0, 4, 8\}$.
\item \textbf{Quantification of Closed-Set Domain Collapse:} We demonstrate that the frozen MobileNetV2 backbone undergoes a 28.56 percentage-point drop in top-1 accuracy---collapsing from 96.92\% [94.71\%, 98.68\%] on the laboratory test partition (Partition A) to 68.36\% [67.20\%, 69.40\%] on pooled external retail fruit (Partition B), with balanced accuracy decreasing from 97.09\% to 68.66\%. Through source stratification, we show that performance varies across external repositories (FruitVision B: 69.14\% accuracy, 69.56\% balanced accuracy; Sultana B: 65.63\% accuracy, 65.55\% balanced accuracy), with misclassifications heavily concentrated in russeted fresh apples and pale citrus cultivars.
\item \textbf{Source-Matched Isolation of Semantic vs. Treatment Shifts:} By benchmarking within shared repositories, source matching reduces cross-repository confounding while potential within-repository capture variation remains. We show that while Energy-based detection achieves robust discrimination under Within-FruitVision semantic shift (Cohort A: 81.57\% AUROC [80.81\%, 82.38\%], 13.18\% operational UFAR), discrimination degrades significantly on external cultivars with visual overlap (Within-Sultana Cohort D: 67.52\% AUROC [65.67\%, 69.25\%], 34.90\% operational UFAR).
\item \textbf{Empirical Diagnosis of Treatment Shift Failure and Operational Trade-Offs:} We demonstrate that standard RGB post-hoc OOD detectors fail to reliably distinguish produce documented under chemical immersion protocols from natural known fruit (Within-FruitVision Cohort B: Energy AUROC of only 53.13\% [52.16\%, 54.09\%]). Under locked validation operating thresholds ($\tau_{95}$), the Dual-Gate cascade yields an operational UFAR-Fresh of 42.83\% [41.26\%, 44.34\%] ($821/1,917$ images) on the treatment-labelled cohort, demonstrating that RGB image features alone cannot reliably identify this condition. Furthermore, we reveal that maintaining an operational UFAR of 16.07\% on external produce incurs an in-distribution known-class false rejection rate (KRR) of approximately 41.98\% ($2,246/5,350$ images).
\end{enumerate}

